% =============================================================================
% COPIONE — Computational Imaging: Deblur & Denoise LBC
% Guida allo studio per l'esposizione orale
% Stile: minimal tech
% =============================================================================
\documentclass[italian,11pt,a4paper]{article}

% --- Minimal tech preamble ---
\usepackage[italian]{babel}
\usepackage[utf8]{inputenc}
\usepackage[T1]{fontenc}
\usepackage{lmodern}
\usepackage{amsmath,amssymb}
\usepackage{geometry}
\geometry{margin=2cm,top=1.8cm,bottom=1.8cm}
\usepackage{booktabs}
\usepackage{multirow}
\usepackage{xcolor}
\usepackage{hyperref}
\usepackage{microtype}
\usepackage{framed}
\usepackage{titlesec}
\usepackage{enumitem}

% --- Palette (tech-minimal) ---
\definecolor{fg}{HTML}{1a1a2e}     % text
\definecolor{accent}{HTML}{00509e}  % deep blue accent
\definecolor{muted}{HTML}{6c757d}   % secondary text
\definecolor{border}{HTML}{dee2e6}  % light borders
\definecolor{codebg}{HTML}{f8f9fa}  % code background

% --- Title style ---
\titleformat{\section}{\Large\bfseries\color{accent}}{}{0em}{}[\vspace{-0.3cm}\rule{\textwidth}{0.5pt}]
\titleformat{\subsection}{\large\bfseries\color{fg}}{}{0em}{}
\titleformat{\subsubsection}{\small\bfseries\color{accent}}{}{0em}{}

\setlist[itemize]{leftmargin=1.2em,noitemsep,topsep=0.15em,label=\textcolor{accent}{\tiny\textbullet}}
\setlist[enumerate]{leftmargin=1.5em,noitemsep,topsep=0.15em}

% --- Custom environments ---
\newenvironment{script}{%
  \smallskip
  \noindent{\small\bfseries\color{accent}$\blacktriangleright$ Da dire}\par
  \vspace{0.05cm}
  \normalsize
}{\par\medskip}

\newenvironment{study}{%
  \smallskip
  \noindent{\small\bfseries\color{muted}$\blacktriangleright$ Da studiare}\par
  \vspace{0.05cm}
  \small\color{fg}
}{\par\medskip}

\newenvironment{qa}{%
  \smallskip
  \noindent{\small\bfseries\color{accent}$\blacktriangleright$ Domande tipiche}\par
  \vspace{0.05cm}
  \small
}{\par\medskip}

\newenvironment{hint}{%
  \smallskip
  \noindent{\small\bfseries\color{teal!70!black}$\blacktriangleright$ Hint + Esempio}\par
  \vspace{0.05cm}
  \small\itshape
}{\par\medskip}

% --- Title ---
\title{\vspace{-1cm}\textcolor{accent}{\textbf{Computational Imaging}} \\[0.1cm]
       Deblur \& Denoise di Immagini di Citologia Cervicale LBC \\[0.2cm]
       \large\textcolor{muted}{Guida allo studio per l'esposizione orale}}
\author{Gruppo R}
\date{A.A. 2025/2026}

\begin{document}
\maketitle
\thispagestyle{empty}
\vspace{0.3cm}
\begin{center}
\small\textcolor{muted}{Accompagna \texttt{slides/presentazione.tex} -- 27 slide, 3 metodi}\\[0.1cm]
\textcolor{muted}{Stampa e usa come traccia durante l'esposizione.}
\end{center}

\newpage
\tableofcontents
\newpage

% =============================================================================
% SLIDE 1
% =============================================================================
\section{Slide 1 -- Title Page}\label{s:01}

\subsection{Contenuto}
Slide iniziale con titolo, autori, affiliazione e link GitHub.

\begin{script}
Buongiorno, siamo il Gruppo R. Presentiamo il nostro progetto di Computational
Imaging: deblur e denoise di immagini di citologia cervicale LBC.

Il progetto confronta tre famiglie metodologiche: Total Variation (variazionale),
UNet (deep learning), DiffPIR (generativo). L'obiettivo e confrontare criticamente
i tre approcci nelle stesse condizioni sperimentali.

Il codice e su GitHub all'indirizzo mostrato in slide.
\end{script}

\begin{study}
\begin{itemize}
\item Tre metodi $\to$ tre famiglie (variazionale, DL, generativo)
\item Confronto equo: stessa pipeline, stesso dataset, stesse metriche
\item Riproducibilita: seed 42, codice open-source
\end{itemize}
\end{study}

\begin{qa}
\textbf{D:} Perche Weighted TV e escluso? \\
\textbf{R:} Siamo in 2, 3 metodi su 4 coprono le famiglie principali. Weighted TV e ibrido.
\end{qa}

% =============================================================================
% SLIDE 2
% =============================================================================
\section{Slide 2 -- Problema Inverso}\label{s:02}

\subsection{Contenuto}
Modello di degradazione $y = Hx + n$, ill-posedness, necessita di regolarizzazione.

\begin{script}
Partiamo dal modello: data un'immagine $x$ (Ground Truth), applichiamo blur
gaussiano $H$ ($\sigma=2$, kernel $9\times9$) e rumore AWGN $n$ a 4 livelli
($\sigma_n = 0.005$--$0.1$). Otteniamo $y$, l'osservazione.

Il problema inverso -- stimare $x$ da $y$ -- e \emph{mal posto} secondo Hadamard:
esistono infinite $x$ compatibili con la stessa $y$, perche il blur cancella le
alte frequenze. Invertire direttamente il blur amplificherebbe il rumore.

Servono \textbf{regolarizzazioni}: la TV usa un prior analitico di
piecewise-smoothness; UNet apprende un prior dai dati; DiffPIR usa un prior
generativo (DDPM).
\end{script}

\begin{study}
\begin{itemize}
\item \textbf{Formula base:} $y = Hx + n$ -- memorizzarla e spiegare ogni termine
\item \textbf{Perche ill-posed:} $H$ cancella info in alta frequenza $\Rightarrow$ soluzione non unica
\item \textbf{Collegamento:} questa slide e il fondamento logico di tutto il progetto
\end{itemize}
\end{study}

\begin{qa}
\textbf{D:} Cosa significa ``mal posto''? \\
\textbf{R:} Hadamard: esiste soluzione? Unica? Dipende con continuita dai dati? Nessuna delle tre e garantita.

\textbf{D:} Perche non possiamo invertire $H$? \\
\textbf{R:} $H$ e mal condizionato: $H^{-1}$ amplifica il rumore in modo catastrofico.
\end{qa}

\begin{hint}
\textbf{Hint:} La convoluzione $Hx$ diventa una moltiplicazione nel dominio di
Fourier. Ecco perche possiamo risolvere il data-fidelity step in FFT con
un semplice prodotto. Invertire nel dominio spaziale richiederebbe una costosa
deconvoluzione.

\textbf{Esempio:} Immagina di guardare attraverso un vetro smerigliato. Il
vetro e il blur $H$: vedi le forme ma non i dettagli. Non puoi ``levare il
vetro'' (invertire $H$) senza amplificare ogni granello di polvere (rumore $n$).
La regolarizzazione e come pulire il vetro senza graffiarlo.
\end{hint}

% =============================================================================
% SLIDE 3
% =============================================================================
\section{Slide 3 -- Metodi Scelti}\label{s:03}

\subsection{Contenuto}
Panoramica dei tre metodi con classificazione per famiglia.

\begin{script}
Tre metodi, tre filosofie diverse.

\textbf{Total Variation} -- metodo variazionale classico. Minimizza
$\|Hx-y\|^2 + \lambda\|\nabla x\|_1$. Non richiede training, e interpretabile,
ma qualita limitata ad alto rumore.

\textbf{UNet} -- deep learning end-to-end. Impara $f_\theta(y) \approx x$ da
coppie (degradata, GT). Inferenza velocissima ($\sim$0.03 s/img), ma richiede
training supervisionato.

\textbf{DiffPIR} -- generativo plug-and-play. Usa un DDPM come prior, alterna
denoising (DDPM) e data-fidelity (risolta in FFT). Produce texture realistiche
ma e lento ($\sim$3 s/img) e allucina a basso rumore.
\end{script}

\begin{study}
\begin{itemize}
\item Saper classificare ogni metodo nella sua famiglia
\item Memorizzare un pro e un contro per ciascuno
\item Tre famiglie diverse $\to$ confronto metodologico
\end{itemize}
\end{study}

% =============================================================================
% SLIDE 4
% =============================================================================
\section{Slide 4 -- Mendeley LBC Cervical Cancer}\label{s:04}

\subsection{Contenuto}
Dataset: 962 immagini, 4 classi diagnostiche, resize $2048\to256$.

\begin{script}
Il dataset Mendeley LBC Cervical Cancer ha 962 immagini di citologia cervicale
in fase liquida, classificate in 4 categorie: NILM (normali, 64\%), HSIL, LSIL e
SCC. Originali a 2048$\times$1536, ridimensionate a 256$\times$256.

Lo sbilanciamento delle classi non e un problema per il restauro: non
classifichiamo, e la varieta morfologica e presente in tutte le classi.
\end{script}

\begin{study}
\begin{itemize}
\item Dataset pubblico, 962 immagini, 4 classi
\item Resize 256: compromesso costo/qualita
\item Per il restauro, lo sbilanciamento NON e un problema
\end{itemize}
\end{study}

% =============================================================================
% SLIDE 5
% =============================================================================
\section{Slide 5 -- Preprocessing e Split}\label{s:05}

\subsection{Contenuto}
Pipeline: resize $\to$ toTensor $\to$ normalize $[-1,1]$ $\to$ split 70/15/15.

\begin{script}
Pipeline condivisa: resize bilineare a 256, conversione in float $[0,1]$,
normalizzazione in $[-1,1]$ (range richiesto dal DDPM, coerente con MSE loss),
split stratificato 70\% train / 15\% val / 15\% test con seed 42.

Tutti i metodi usano lo stesso test set: 145 immagini.
\end{script}

\begin{study}
\begin{itemize}
\item Perche $[-1,1]$? DDPM addestrato con rumore gaussiano centrato
\item Split stratificato: mantiene le proporzioni delle classi
\item Seed 42: riproducibilita
\end{itemize}
\end{study}

% =============================================================================
% SLIDE 6
% =============================================================================
\section{Slide 6 -- Degradazione}\label{s:06}

\subsection{Contenuto}
Pipeline di degradazione: blur gaussiano + AWGN a 4 livelli.

\begin{script}
Degradazione identica per tutti i metodi: blur gaussiano $\sigma=2$ kernel
$9\times9$, poi AWGN a 4 livelli: $\sigma_n = 0.005,\;0.01,\;0.05,\;0.1$.

I 4 livelli coprono regimi qualitativamente distinti: da rumore trascurabile
a rumore pesante. Questo permette di studiare la transizione.
\end{script}

\begin{study}
\begin{itemize}
\item Pipeline \emph{identica} per tutti -- confronto equo
\item I 4 livelli non sono arbitrari: ognuno rappresenta un regime
\end{itemize}
\end{study}

% =============================================================================
% SLIDE 7
% =============================================================================
\section{Slide 7 -- Total Variation: Teoria}\label{s:07}

\subsection{Contenuto}
Formulazione TV: $\hat{x} = \arg\min\|Hx-y\|^2 + \lambda\cdot TV(x)$. $TV(x) = \|\nabla x\|_1$.

\begin{script}
La TV minimizza un funzionale a due termini: data-fidelity ($\|Hx-y\|^2$) e
regolarizzazione ($\lambda\cdot\|\nabla x\|_1$).

La TV e la norma L1 del gradiente. Usare L1 e la scelta fondamentale: preserva
i bordi. Se usassimo L2 (Tikhonov), i bordi verrebbero sfumati.

$\lambda$ controlla il bilanciamento. Il problema e convesso -- soluzione
globalmente ottima.
\end{script}

\begin{study}
\begin{itemize}
\item \textbf{L1 $\to$ preserva bordi; L2 $\to$ sfuma bordi}
\item Staircasing: artefatti a gradino se $\lambda$ troppo alto
\item Convessita: minimo globale garantito
\end{itemize}
\end{study}

\begin{qa}
\textbf{D:} Perche L1 e non L2? \\
\textbf{R:} L1 promuove sparsity nel gradiente (pochi bordi netti), L2 produce
soluzioni dense (tante piccole variazioni, bordi sfumati).
\end{qa}

\begin{hint}
\textbf{Hint:} L1 e L2 penalizzano diversamente. L1 somma i valori assoluti,
L2 somma i quadrati. Un gradiente grande (bordo) pesa tanto in L1, ma
\emph{moltissimo} in L2. Un gradiente piccolo (sfumatura) pesa poco in L2.
Risultato: L2 preferisce tante piccole variazioni (sfumatura), L1 preferisce
pochi grandi salti (bordi netti).

\textbf{Esempio:} Un gradino alto 1 su 1 pixel: L1 = 1, L2 = 1. Lo stesso
gradino sfumato su 10 pixel (0.1 ciascuno): L1 = $10 \times 0.1 = 1$ (uguale),
L2 = $10 \times 0.01 = 0.1$ (10 volte meno!). L2 ``pensa'' che la versione
sfumata sia migliore, perche penalizza meno -- ma noi vogliamo il bordo netto!
\end{hint}

% =============================================================================
% SLIDE 8
% =============================================================================
\section{Slide 8 -- Total Variation: Algoritmo}\label{s:08}

\subsection{Contenuto}
Ottimizzazione con Adam, 150 iter, $\lambda=0.005$, clamping $[-1,1]$.

\begin{script}
Implementazione pratica: inizializziamo $x=y$, usiamo Adam (lr 0.001), 150
iterazioni, $\lambda=0.005$. Ad ogni passo: blur, loss, backprop, update,
clamping in $[-1,1]$.

Perche Adam? I gradienti TV hanno magnitudo variabile (grandi sui bordi, piccoli
nelle zone uniformi). Adam adatta il learning rate.

150 iter bastano per convergenza empirica. $\lambda=0.005$ e il miglior
compromesso.
\end{script}

\begin{study}
\begin{itemize}
\item Adam: learning rate adattivo, meglio di GD classico
\item 150 iter + clamping: trade-off velocita/qualita
\item $\lambda=0.005$: bilanciamento ottimale
\end{itemize}
\end{study}

\begin{hint}
\textbf{Hint:} Adam combina momentum (accelerazione nelle direzioni
consistenti) con adaptive learning rate (passi piccoli dove il gradiente e
variabile). Per la TV, il gradiente e grande ai bordi (discontinuita) e
piccolo nelle zone uniformi -- Adam si adatta automaticamente. SGD classico
dovrebbe usare un learning rate unico per tutti i pixel.

\textbf{Esempio:} E come un taxi in citta: in autostrada (zona uniforme,
gradiente piccolo) accelera, nel traffico (bordo, gradiente grande) frena.
SGD invece va sempre a 50 km/h, ovunque.
\end{hint}

% =============================================================================
% SLIDE 9
% =============================================================================
\section{Slide 9 -- Total Variation: Risultati}\label{s:09}

\subsection{Contenuto}
PSNR 32.09 dB ($\sigma_n=0.005$) $\to$ 26.54 dB ($\sigma_n=0.1$). SSIM 0.911 $\to$ 0.586.

\begin{script}
TV eccelle a basso rumore: PSNR $>$ 32 dB, SSIM $>$ 0.9 a $\sigma_n \leq 0.01$,
senza training. Ad alto rumore, degrada: staircasing diventa dominante, PSNR
scende a 26.54 dB, SSIM a 0.586.

La TV funziona bene quando l'assunzione piecewise-smooth e valida -- cioe a
basso rumore.
\end{script}

\begin{study}
\begin{itemize}
\item Eccellente a basso rumore (PSNR $>$ 32 dB)
\item Degrada ad alto rumore (staircasing)
\item Vantaggio: zero training richiesto
\end{itemize}
\end{study}

% =============================================================================
% SLIDE 10
% =============================================================================
\section{Slide 10 -- UNet: Architettura}\label{s:10}

\subsection{Contenuto}
Encoder-decoder con skip connections, GroupNorm, input 4 canali, $\sim$1.9M params.

\begin{script}
UNet: encoder comprime (16$\to$256 canali), decoder ricostruisce. Skip
connections preservano dettagli spaziali. Ogni blocco: DoubleConv (Conv3x3 +
GroupNorm + ReLU $\times$2).

GroupNorm: piu stabile di BatchNorm con batch size piccolo (16). Input: 4 canali
(RGB + noise map). Output: RGB + tanh. $\sim$1.9M parametri -- snello, gira su CPU.
\end{script}

\begin{study}
\begin{itemize}
\item Skip connections: preservano informazione spaziale
\item GroupNorm $>$ BatchNorm per batch piccoli
\item 4 canali input: condizionamento sul rumore
\end{itemize}
\end{study}

\begin{hint}
\textbf{Hint:} Le skip connections collegano l'encoder (dove l'immagine e
grande e dettagliata) al decoder (dove l'immagine viene ricostruita). Senza,
il collo di bottiglia (bottleneck) perderebbe la posizione dei dettagli --
sapresti \emph{cosa} c'e ma non \emph{dove} si trova.

GroupNorm: divide i canali in gruppi e normalizza ogni gruppo. BatchNorm
normalizza su tutto il batch -- se il batch e piccolo, la stima e rumorosa.
GroupNorm funziona bene anche con batch = 1.

\textbf{Esempio:} Le skip connections sono come un pittore che tiene la foto
originale a fianco mentre dipinge: guarda il soggetto (encoder), poi dipinge
(decoder), ma ogni tanto torna a controllare i dettagli originali (skip).
\end{hint}

% =============================================================================
% SLIDE 11
% =============================================================================
\section{Slide 11 -- UNet: Training}\label{s:11}

\subsection{Contenuto}
Loss L1, Adam, batch 16, 50 epoche, multi-noise augmentation.

\begin{script}
Loss L1 (piu robusta della MSE agli outlier -- i dettagli fini). Adam lr
$10^{-4}$, batch 16, 50 epoche su 673 immagini.

Multi-noise augmentation: ad ogni batch, $\sigma_n$ scelto uniformemente tra i 4
livelli. Il modello impara a gestire tutto lo spettro di rumore.
\end{script}

\begin{study}
\begin{itemize}
\item L1 $>$ MSE: preserva dettagli fini
\item Multi-noise: chiave della robustezza
\item Training su CPU: 50 epoche, poche ore
\end{itemize}
\end{study}

% =============================================================================
% SLIDE 12
% =============================================================================
\section{Slide 12 -- UNet: Risultati}\label{s:12}

\subsection{Contenuto}
PSNR 28.46--29.79 dB (escursione $\sim$1 dB). Inferenza 0.03 s/img.

\begin{script}
Due risultati notevoli: degradazione trascurabile col rumore ($\sim$1 dB di
escursione su tutto lo spettro), e inferenza velocissima (0.03 s/img -- 200$\times$
piu veloce di TV, 100$\times$ di DiffPIR).

A $\sigma_n=0.1$, UNet supera TV di quasi 2 dB. A basso rumore, TV e ancora
leggermente superiore: UNet generalizza ma paga un piccolo costo.
\end{script}

\begin{study}
\begin{itemize}
\item Minima degradazione PSNR ($\sim$1 dB su tutto lo spettro)
\item Inferenza 0.03 s/img: 200$\times$ piu veloce di TV
\item Supera TV a $\sigma_n=0.1$
\end{itemize}
\end{study}

% =============================================================================
% SLIDE 13
% =============================================================================
\section{Slide 13 -- DiffPIR: Panoramica}\label{s:13}

\subsection{Contenuto}
DiffPIR (CVPR 2023) = DDPM + Plug-and-Play. LightUNet custom (1.26M params).

\begin{script}
DiffPIR integra un DDPM in un framework Plug-and-Play: alterna denoising
(DDPM) e data-fidelity (Wiener in FFT).

La novita: invece di usare un modello pre-addestrato su ImageNet (500M params,
2 GB), abbiamo addestrato una LightUNet custom da zero su 673 immagini LBC.
1.26M parametri, $\sim$5 MB, gira su CPU.
\end{script}

\begin{study}
\begin{itemize}
\item Plug-and-Play: denoising e data-fidelity separati
\item Fidelity risolta in FFT: soluzione analitica
\item LightUNet custom: specifica per citologia
\end{itemize}
\end{study}

% =============================================================================
% SLIDE 14
% =============================================================================
\section{Slide 14 -- DDPM e LightUNet}\label{s:14}

\subsection{Contenuto}
DDPM: forward $q(x_t|x_{t-1})$, reverse $p_\theta(x_{t-1}|x_t)$, MSE su $\varepsilon$.

\begin{script}
DDPM: forward aggiunge rumore in 1000 step (scheduler lineare
$\beta_1=10^{-4}$, $\beta_T=0.02$). Reverse: il modello predice il rumore
$\varepsilon$ aggiunto, non l'immagine.

LightUNet: canali base 32, encoder 32$\to$64$\to$128, time embedding sinusoidale
+ MLP per il timestep, GroupNorm + SiLU. Addestrata 30 epoche su 673 immagini.
\end{script}

\begin{study}
\begin{itemize}
\item Il modello predice $\varepsilon$, non $x_0$
\item Time embedding: codifica il timestep corrente
\item 1000 step forward, ma sampling con 15 step DDIM
\end{itemize}
\end{study}

\begin{hint}
\textbf{Hint:} Perche il DDPM predice $\varepsilon$ (rumore) e non $x_0$
(immagine pulita)? Perche $\varepsilon$ ha una distribuzione semplice e nota
(gaussiana standard), mentre $x_0$ ha una distribuzione complicatissima
(immagini reali). E piu facile imparare a ``togliere lo sporco'' che a
``disegnare da zero''.

Time embedding: il modello riceve il timestep $t$ codificato con seni/coseni
(come le coordinate in NeRF/Transformer). Cosi sa se deve fare un passo
piccolo ($t$ piccolo, immagine quasi pulita) o grande ($t$ grande, molto
rumore).

\textbf{Esempio:} E come insegnare a qualcuno a pulire una vecchia foto:
invece di ridisegnarla, impara a riconoscere e rimuovere graffi e macchie
($\varepsilon$). Il time embedding e come dirgli ``questa foto e molto sporca
(fai un passo grosso)'' o ``e quasi pulita (sfuma solo l'ultimo graffio)''.
\end{hint}

% =============================================================================
% SLIDE 15
% =============================================================================
\section{Slide 15 -- Algoritmo DiffPIR}\label{s:15}

\subsection{Contenuto}
Passi: denoising $\to$ data-fidelity (FFT) $\to$ DDIM step. $\rho_t$ dinamico.

\begin{script}
Ad ogni iterazione: (1) denoising -- DDPM stima $\hat{x}_0$ da $x_t$;
(2) data-fidelity -- soluzione di Wiener in FFT: bilancia $\hat{x}_0$ con
l'osservazione $y$; (3) DDIM step -- sample $x_{t-1}$ da $\tilde{x}_0$.

$\rho_t$ e dinamico: alto a $t$ grande (prior pesa di piu), basso a $t$ piccolo
(pesano i dati). Usiamo 15 step DDIM anziche 1000 per accelerare.
\end{script}

\begin{study}
\begin{itemize}
\item Tre passi: denoise, fidelity, sample
\item Fidelity in FFT: convoluzione $\to$ moltiplicazione (efficiente)
\item $\rho_t$ dinamico: bilanciamento adattivo
\item 15 step DDIM: sampling sub-campionato
\end{itemize}
\end{study}

\begin{hint}
\textbf{Hint:} $\rho_t$ e un parametro che bilancia prior e dati. All'inizio
($t$ grande, tanto rumore) $\rho_t$ e alto -- ci fidiamo del prior perche
l'osservazione e troppo rumorosa. Alla fine ($t$ piccolo, poco rumore) $\rho_t$
e basso -- ci fidiamo dei dati perche sono piu affidabili del prior.

La soluzione di Wiener in FFT e una formula chiusa: $\hat{x}_0 = \frac{\dots}$,
calcolabile con una FFT forward e una inverse. Senza FFT, dovremmo risolvere
un sistema lineare $O(N^3)$; in FFT e $O(N \log N)$.

\textbf{Esempio:} $\rho_t$ e come un arbitro di calcio: all'inizio della partita
($t$ grande, caos) lascia giocare (prior), alla fine ($t$ piccolo) fischia di
piu (dati). Il passaggio graduale evita strappi.
\end{hint}

% =============================================================================
% SLIDE 16
% =============================================================================
\section{Slide 16 -- DiffPIR: Risultati}\label{s:16}

\subsection{Contenuto}
PSNR 15.78 dB ($\sigma_n=0.005$) $\to$ 25.46 dB ($\sigma_n=0.1$).

\begin{script}
DiffPIR ha comportamento inverso: PSNR \emph{cresce} col rumore. Da 15.78 dB a
$\sigma_n=0.005$ (peggiore di tutti) a 25.46 dB a $\sigma_n=0.1$ (miglioramento).

A basso rumore: il modello \emph{allucina} dettagli -- il prior generativo
sovrascrive quelli reali. Ad alto rumore: il prior aiuta a ricostruire cio che
manca.

Questo e un esempio perfetto di \textbf{bias-variance trade-off}: a $\sigma_n$
basso il bias domina, a $\sigma_n$ alto la varianza e ridotta dal prior.
\end{script}

\begin{study}
\begin{itemize}
\item \textbf{Trend inverso:} PSNR migliora col rumore
\item Allucinazione a basso $\sigma_n$: prior sovrascrive dettagli
\item Bias-variance trade-off: menzionarlo all'esame!
\end{itemize}
\end{study}

\begin{hint}
\textbf{Hint:} Bias-variance trade-off e il cuore del machine learning. Un
modello con \textbf{alto bias} fa assunzioni forti (es. ``l'immagine e
piecewise-smooth'' come TV) -- sbaglia sistematicamente ma e prevedibile.
Un modello con \textbf{alta varianza} e molto flessibile (es. una rete
profonda) -- azzecca quando ha dati buoni, ma e instabile con dati poveri.

DiffPIR ha forte prior generativo (alto bias). A $\sigma_n$ basso questo bias
e dannoso: il modello ``crede'' che l'immagine debba assomigliare a quelle
viste in training e sovrascrive i dettagli reali. A $\sigma_n$ alto il rumore
rende i dati inaffidabili, e avere un prior forte \emph{aiuta} a stabilizzare.

\textbf{Esempio:} Due fotografi a un concerto. Il primo (alto bias) sa suonare
solo rock -- se il concerto e rock, foto perfette; se e jazz, sbaglia tutto.
Il secondo (alta varianza) fotografa qualsiasi genere -- ma se c'e poca luce
(molto rumore), le foto vengono mosse. DiffPIR e il primo fotografo: ottimo
se conosce gia il genere, male se e sorpreso.
\end{hint}

% =============================================================================
% SLIDE 17
% =============================================================================
\section{Slide 17 -- Implementazione}\label{s:17}

\subsection{Contenuto}
Struttura repository: moduli separati, pipeline condivisa, 34 test.

\begin{script}
Repo modulare: ogni metodo in \texttt{src/methods/} con script indipendenti.
Pipeline condivisa: degradazione, metriche, plotting. 34 test unitari. Seed 42
per riproducibilita.
\end{script}

\begin{study}
\begin{itemize}
\item Modularita: ogni metodo e indipendente
\item Pipeline condivisa: confronto equo
\item 34 test: copertura di degradazione, metriche, metodi
\end{itemize}
\end{study}

% =============================================================================
% SLIDE 18
% =============================================================================
\section{Slide 18 -- Confronto PSNR / SSIM}\label{s:18}

\subsection{Contenuto}
Tabella riassuntiva dei risultati quantitativi.

\begin{script}
Tabella centrale del confronto. TV domina a $\sigma_n \leq 0.01$ (PSNR $>$ 32 dB,
SSIM $>$ 0.9). A $\sigma_n=0.05$: TV ha PSNR piu alto, UNet ha SSIM migliore
(0.864 vs 0.837). A $\sigma_n=0.1$: UNet nettamente superiore (PSNR 28.46 vs 26.54,
SSIM 0.795 vs 0.586). DiffPIR migliora col rumore ma resta sotto.
\end{script}

\begin{study}
\begin{itemize}
\item TV: eccellente a basso rumore
\item UNet: robusto su tutto lo spettro, migliore ad alto
\item DiffPIR: complementare, migliora col rumore
\end{itemize}
\end{study}

\begin{hint}
\textbf{Hint:} PSNR e SSIM misurano cose diverse. PSNR = errore quadratico
medio in dB: confronta pixel per pixel. SSIM = similarita strutturale: guarda
luminanza, contrasto e struttura in finestre locali. Per questo TV ha PSNR
decente ma SSIM basso ad alto rumore: lo staircasing sposta i pixel di poco
(PSNR regge) ma distrugge la struttura locale (SSIM crolla).

\textbf{Esempio:} PSNR e come confrontare due foto misurando la distanza
pixel-per-pixel con un righello. SSIM e come chiedere ``queste due foto
sembrano la stessa scena?''. La prima e precisa ma ingenua, la seconda coglie
la qualita percepita.
\end{hint}

% =============================================================================
% SLIDE 19
% =============================================================================
\section{Slide 19 -- Grafico Comparativo}\label{s:19}

\subsection{Contenuto}
Grafico PSNR e SSIM vs $\sigma_n$ per i tre metodi.

\begin{script}
Il grafico visualizza le metriche. Linee TV e UNet quasi piatte in PSNR,
divergenti in SSIM -- perche il PSNR e meno sensibile alla struttura. TV crolla
in SSIM dopo $\sigma_n=0.05$ (staircasing). DiffPIR ha pendenza crescente.

L'incrocio TV/UNet a $\sigma_n\approx0.05$ segna il punto in cui il metodo
classico cede al deep learning.
\end{script}

\begin{study}
\begin{itemize}
\item PSNR piatto, SSIM divergente: perche? Struttura vs pixel
\item Incrocio TV/UNet a $\sigma_n\approx0.05$
\item DiffPIR: trend crescente in entrambe le metriche
\end{itemize}
\end{study}

% =============================================================================
% SLIDE 20
% =============================================================================
\section{Slide 20 -- Griglia Completa}\label{s:20}

\subsection{Contenuto}
Griglia visiva: metodi $\times$ livelli di rumore.

\begin{script}
Griglia che confronta visivamente i metodi per ogni livello di rumore.
TV: pulita a sinistra, degrada a destra. UNet: uniforme su tutta la griglia.
DiffPIR: migliora da sinistra a destra. Conferma visiva delle metriche.
\end{script}

\begin{study}
\begin{itemize}
\item Correlazione visiva-metriche: la griglia conferma la tabella
\end{itemize}
\end{study}

% =============================================================================
% SLIDE 21
% =============================================================================
\section{Slide 21 -- Confronto Visivo $\sigma_n=0.01$}\label{s:21}

\subsection{Contenuto}
Confronto visivo a basso rumore.

\begin{script}
A $\sigma_n=0.01$, TV e UNet sono eccellenti, quasi indistinguibili dalla GT.
DiffPIR mostra gia artefatti generativi: allucina dettagli inesistenti.
Il crop ingrandito evidenzia le differenze nei nuclei cellulari.
\end{script}

\begin{study}
\begin{itemize}
\item A basso rumore: TV e UNet funzionano, DiffPIR no
\item Crop: dettaglio sui nuclei
\end{itemize}
\end{study}

% =============================================================================
% SLIDE 22
% =============================================================================
\section{Slide 22 -- Confronto Visivo $\sigma_n=0.05$}\label{s:22}

\subsection{Contenuto}
Confronto visivo a medio rumore.

\begin{script}
A $\sigma_n=0.05$, primo cambio di regime. TV mostra staircasing: l'immagine
e squadrata. UNet mantiene buona qualita -- inizia il vantaggio del DL. DiffPIR
migliora rispetto a $\sigma_n=0.01$: texture piu plausibili.
\end{script}

\begin{study}
\begin{itemize}
\item TV: staircasing visibile
\item UNet: inizia a superare TV (visivamente)
\item DiffPIR: texture piu realistiche
\end{itemize}
\end{study}

% =============================================================================
% SLIDE 23
% =============================================================================
\section{Slide 23 -- Confronto Visivo $\sigma_n=0.1$}\label{s:23}

\subsection{Contenuto}
Confronto visivo ad alto rumore + mappe di differenza.

\begin{script}
A $\sigma_n=0.1$, le differenze sono massime. TV: staircasing dominante, PSNR
26.54 dB, SSIM 0.586. UNet: il migliore, PSNR 28.46 dB, pulito e riconoscibile.
DiffPIR: texture realistiche ma PSNR 25.46 dB.

Le mappe di differenza ($|\hat{x} - x_{\text{GT}}|$) sono rivelatrici: TV ha
errori strutturati a blocchi, UNet ha errori minimi e isotropi, DiffPIR ha
errori concentrati dove ha allucinato.
\end{script}

\begin{study}
\begin{itemize}
\item $\sigma_n=0.1$: massima divergenza tra i metodi
\item Mappe di differenza: errore strutturato (TV) vs isotropo (UNet)
\end{itemize}
\end{study}

% =============================================================================
% SLIDE 24
% =============================================================================
\section{Slide 24 -- Confronto tra le Famiglie}\label{s:24}

\subsection{Contenuto}
Pro e contro di ogni famiglia metodologica.

\begin{script}
Riassunto: TV -- no training, interpretabile, eccellente a basso rumore ma
staircasing. UNet -- velocissimo ($0.03$ s/img), robusto, ma richiede dati
supervisionati. DiffPIR -- texture realistiche, migliora col rumore, ma lento
($3$ s/img) e allucina a basso rumore.

Nessuno domina tutti i regimi. La scelta dipende dal contesto.
\end{script}

\begin{study}
\begin{itemize}
\item Tre trade-off diversi: velocita, qualita, generalizzazione
\item Nessun metodo universalmente migliore
\end{itemize}
\end{study}

% =============================================================================
% SLIDE 25
% =============================================================================
\section{Slide 25 -- Regimi di Funzionamento}\label{s:25}

\subsection{Contenuto}
Tabella scenario $\times$ metodo: chi scegliere in ogni regime.

\begin{script}
Tabella riassuntiva: basso rumore $\to$ TV (eccellente); medio rumore $\to$
UNet (molto buono); alto rumore $\to$ UNet (migliore), DiffPIR (buono).
Velocita: UNet ($0.03$ s) $>$ DiffPIR ($3$ s) $>$ TV ($7$ s).
\end{script}

\begin{study}
\begin{itemize}
\item Basso $\sigma_n$: TV
\item Medio-alto $\sigma_n$: UNet
\item Alto $\sigma_n$: considerare DiffPIR se texture importanti
\end{itemize}
\end{study}

% =============================================================================
% SLIDE 26
% =============================================================================
\section{Slide 26 -- Conclusioni}\label{s:26}

\subsection{Contenuto}
Conclusioni e takeaway del progetto.

\begin{script}
Tre risultati principali:
(1) UNet e il miglior compromesso: robusto su tutto lo spettro, velocissimo.
(2) Il confronto equo e il vero contributo: stessa pipeline, stesse metriche.
(3) Nessun metodo domina tutti i regimi: la scelta e contesto-dipendente.

Grazie per l'attenzione!
\end{script}

\begin{study}
\begin{itemize}
\item Tre takeaway chiari
\item Enfatizzare il confronto equo come contributo
\end{itemize}
\end{study}

% =============================================================================
% SLIDE 27
% =============================================================================
\section{Slide 27 -- Riferimenti}\label{s:27}

\subsection{Contenuto}
Bibliografia: DiffPIR (CVPR 2023), DDPM (NeurIPS 2020), DDIM (ICLR 2021), etc.

\begin{study}
\begin{itemize}
\item Conoscere i paper principali (almeno titolo e conferenza)
\item DiffPIR: Zhu et al., CVPR 2023
\item DDPM: Ho et al., NeurIPS 2020
\item DDIM: Song et al., ICLR 2021
\end{itemize}
\end{study}

% =============================================================================
% Appendice: strategia di esposizione
% =============================================================================
\newpage
\section*{Strategia di Esposizione}
\addcontentsline{toc}{section}{Strategia di Esposizione}

\subsection*{Prima dell'esame}
\begin{enumerate}
\item Leggi il copione slide per slide
\item Per ogni slide, spiega il contenuto \textbf{a voce alta} senza leggere
\item Verifica di saper rispondere alle domande tipiche
\item Esercitati con un timer: 15--20 minuti totali
\end{enumerate}

\subsection*{Durante l'esposizione}
\begin{itemize}
\item \textbf{Slide 2} (Problema Inverso): parti da $y = Hx + n$ -- e il fondamento
\item \textbf{Slide 7} (TV): enfatizza L1 vs L2 e staircasing
\item \textbf{Slide 18} (Confronto): usa la tabella come punto di riferimento
\item \textbf{Slide 23} (Mappe differenza): mostra visivamente le differenze
\item \textbf{Conclusione}: UNet miglior compromesso, ma non universalmente
\end{itemize}

\subsection*{Collegamenti da sottolineare}
\begin{itemize}
\item \textbf{Ottimizzazione:} norma L1 $\to$ sparsity; Adam vs GD
\item \textbf{Probabilita:} DDPM forward/reverse; bias-variance trade-off
\item \textbf{Elaborazione segnali:} convoluzione $\to$ FFT; Wiener filter
\end{itemize}

\subsection*{Domande frequenti}
\begin{enumerate}
\item ``Qual e il metodo migliore?'' $\to$ Dipende dal regime e dalla metrica
\item ``Perche non avete usato Weighted TV?'' $\to$ 3 su 4 per 2 studenti
\item ``Perche L1 e non L2 in TV?'' $\to$ Preserva bordi, sparsity
\item ``Perche DiffPIR migliora col rumore?'' $\to$ Bias-variance trade-off
\item ``Come garantite riproducibilita?'' $\to$ Seed 42, pipeline unica, test
\end{enumerate}

\end{document}
